\documentclass[14pt, a4paper, landscape]{extarticle}

\usepackage{xcolor}
\usepackage{geometry}
\usepackage{graphicx}
\usepackage{booktabs}
\usepackage{array}
\usepackage{longtable}
\usepackage{titlesec}
\usepackage{enumitem}
\usepackage{tcolorbox}
\usepackage{fancyhdr}
\usepackage{tikz}

\geometry{margin=1cm}

\usepackage{xepersian}
\settextfont{Vazirmatn}
\setlatintextfont{Times New Roman}
\setdigitfont{Vazirmatn}

\definecolor{primary}{RGB}{139, 90, 43}
\definecolor{accent}{RGB}{224, 122, 63}
\definecolor{bg}{RGB}{245, 240, 232}
\definecolor{darkbg}{RGB}{92, 58, 26}

\pagestyle{fancy}
\fancyhf{}
\fancyhead[R]{\small\textcolor{primary}{مهندسی نرم‌افزار — جلسهٔ ۱}}
\fancyfoot[C]{\textcolor{primary}{\thepage}}
\renewcommand{\headrulewidth}{0.5pt}
\setlength{\headheight}{15pt}

\setlength{\parindent}{0pt}
\setlength{\parskip}{0.5em}

\newcommand{\slide}[1]{%
    \newpage
    \begin{center}
    \begin{minipage}[t]{0.95\textwidth}
    \vspace{0.5cm}
    #1
    \vspace{0.5cm}
    \end{minipage}
    \end{center}
}

\newcommand{\slidetitle}[1]{%
    \begin{center}
        {\LARGE\bfseries\textcolor{primary}{#1}}
    \end{center}
    \vspace{0.5cm}
}

\newcommand{\myaccent}[1]{\textcolor{accent}{#1}}

\begin{document}

\pagestyle{empty}

\begin{tikzpicture}[remember picture, overlay]
    \fill[darkbg] (current page.north west) rectangle (current page.south east);
\end{tikzpicture}

\begin{center}
    \vspace*{4cm}
    {\color{white}\Huge\bfseries مهندسی نرم‌افزار}
    
    \vspace{1cm}
    
    {\color{accent}\LARGE جلسهٔ ۱ — معرفی درس}
    
    \vspace{2cm}
    
    {\color{white}\Large ابراهیم فضلی}
    
    \vspace{0.5cm}
    
    {\color{white}\large ترم ۱۴۰۴}
\end{center}

\pagestyle{fancy}

% ---------- اسلاید ۲ ----------
\slide{
    \slidetitle{یک تصویر، یک سؤال}
    \vspace{1.5cm}
    \begin{center}
        {\Large\bfseries نرم‌افزار هم می‌تواند شکست بخورد.}
        
        \vspace{1.5cm}
        
        {\Huge\bfseries\myaccent{چرا؟}}
    \end{center}
}

% ---------- اسلاید ۳ ----------
\slide{
    \slidetitle{آمار تکان‌دهنده}
    \vspace{1cm}
    \begin{center}
        {\fontsize{80}{90}\selectfont\bfseries\myaccent{۷۰٪}}
        
        \vspace{1cm}
        
        {\Large پروژه‌های نرم‌افزاری شکست می‌خورند،}
        
        {\Large دیر تحویل می‌شوند یا از بودجه فراتر می‌روند.}
        
        \vspace{1.5cm}
        
        {\small منبع: Standish Group Chaos Report}
    \end{center}
}

% ---------- اسلاید ۴ ----------
\slide{
    \slidetitle{به نظر شما چرا پروژه‌ها شکست می‌خورند؟}
    \vspace{1cm}
    \begin{center}
        {\Huge\bfseries\textcolor{primary}{نیازمندی‌ها؟ زمان؟ تیم؟ ابزار؟}}
        
        \vspace{1.5cm}
        
        {\large\myaccent{(۳۰ ثانیه فکر کنید، بعد چند نفر بگویند)}}
    \end{center}
}

% ---------- اسلاید ۵ ----------
\slide{
    \slidetitle{کدنویس vs مهندس نرم‌افزار}
    \vspace{0.5cm}
    \begin{center}
        \begin{tabular}{p{6cm}p{6cm}}
            \toprule
            \textbf{\large کدنویس} & \textbf{\large مهندس نرم‌افزار} \\
            \midrule
            کد می‌نویسد & مسئله را حل می‌کند \\
            \hline
            تنها کار می‌کند & در تیم کار می‌کند \\
            \hline
            به امروز فکر می‌کند & به ۵ سال بعد فکر می‌کند \\
            \hline
            تست را ضروری نمی‌داند & تست را بخشی از کار می‌داند \\
            \hline
            مشتری را نمی‌شناسد & مشتری را می‌فهمد \\
            \bottomrule
        \end{tabular}
    \end{center}
}

% ---------- اسلاید ۶ ----------
\slide{
    \begin{center}
        \vspace{2.5cm}
        {\Huge\bfseries مهندسی نرم‌افزار}
        
        \vspace{1.5cm}
        
        {\LARGE\myaccent{= ساختن چیزی که کار می‌کند،}}
        
        \vspace{0.5cm}
        
        {\LARGE\myaccent{در زمان مشخص، با تیم، و قابل نگهداری.}}
    \end{center}
}

% ---------- اسلاید ۷ ----------
\slide{
    \slidetitle{اهداف درس}
    \vspace{0.5cm}
    {\large در پایان این درس می‌توانید:}
    \vspace{0.5cm}
    \begin{enumerate}[itemsep=0.6em]
        \item نیازمندی‌ها را از مشتری استخراج کنید.
        \item معماری و طراحی یک سیستم را انجام دهید.
        \item در یک تیم چابک کار کنید.
        \item کد را تست، بازبینی و CI/CD تحویل دهید.
        \item کیفیت، ریسک و بدهی فنی را ارزیابی کنید.
    \end{enumerate}
}

% ---------- اسلاید ۸ ----------
\slide{
    \slidetitle{نقشهٔ ۱۴ هفته}
    \vspace{0.3cm}
    \begin{center}
        \begin{tabular}{cl}
            \toprule
            \textbf{هفته} & \textbf{موضوع} \\
            \midrule
            ۱-۲ & نیازمندی‌ها \\
            \hline
            ۳-۴ & مدل‌سازی و معماری \\
            \hline
            ۵ & Sprint 1 Review \\
            \hline
            ۶-۷ & طراحی شیءگرا و الگوها \\
            \hline
            ۸ & میان‌ترم \\
            \hline
            ۹-۱۰ & تست و CI/CD \\
            \hline
            ۱۱ & Sprint 2 Review \\
            \hline
            ۱۲-۱۳ & کیفیت، ریسک، DevOps \\
            \hline
            ۱۴ & Demo Day \\
            \bottomrule
        \end{tabular}
    \end{center}
}

% ---------- اسلاید ۹ ----------
\slide{
    \slidetitle{منابع}
    \vspace{0.5cm}
    \begin{itemize}[itemsep=0.6em]
        \item \textbf{اصلی:} Software Engineering – Sommerville
        \item \textbf{تکمیلی:} Clean Code – Robert C. Martin
        \item \textbf{رایگان:} Scrum Guide، The Pragmatic Programmer
        \item لینک به منابع آنلاین در سامانهٔ درس
    \end{itemize}
}

% ---------- اسلاید ۱۰ ----------
\slide{
    \slidetitle{انتظارات}
    \vspace{0.5cm}
    \begin{itemize}[itemsep=0.6em]
        \item حضور فعال در کلاس
        \item مشارکت در کار تیمی
        \item تحویل به‌موقع Sprintها
        \item استفاده از AI مجاز است، اما با شفافیت
        \item کپی = نمرهٔ صفر
    \end{itemize}
}

% ---------- اسلاید ۱۱ ----------
\slide{
    \slidetitle{قاعدهٔ طلایی این درس}
    \vspace{1.5cm}
    \begin{center}
        {\Large هر مفهوم را اول روی پروژهٔ من می‌بینید،}
        
        \vspace{0.5cm}
        
        {\Large بعد همان را روی پروژهٔ خودتان پیاده می‌کنید.}
        
        \vspace{2cm}
        
        {\Large\bfseries\myaccent{من = هوم‌استور}}
        
        \vspace{0.5cm}
        
        {\Large\bfseries\myaccent{شما = پروژهٔ تیم خودتان}}
    \end{center}
}

% ---------- اسلاید ۱۲ ----------
\slide{
    \slidetitle{پروژهٔ من: هوم‌استور}
    \vspace{0.5cm}
    \begin{center}
        {\Large\bfseries یک فروشگاه آنلاین لوازم خانگی}
    \end{center}
    
    \vspace{1cm}
    
    \begin{itemize}[itemsep=0.5em]
        \item من در طول ترم این پروژه را زنده می‌سازم.
        \item هدف: نشان دادن نمونهٔ عملی از هر مفهوم درس.
        \item شما هر هفته پیشرفت آن را در GitHub می‌بینید.
        \item لازم نیست کپی کنید؛ فقط اصول را یاد بگیرید.
    \end{itemize}
}

% ---------- اسلاید ۱۳ ----------
\slide{
    \slidetitle{مشتری من}
    \vspace{1cm}
    \begin{center}
        {\LARGE\bfseries\myaccent{خانم رضایی، ۴۲ ساله}}
        \par\vspace{1.5cm}
        {\Large\itshape «محصولاتم اصل و با گارانتی‌اند،\\
        اما فقط مشتری‌های محلی من را می‌شناسند.»}
    \end{center}
}

% ---------- اسلاید ۱۴ ----------
\slide{
    \slidetitle{مشکل فعلی پروژهٔ من}
    \vspace{0.5cm}
    \begin{itemize}[itemsep=0.6em]
        \item فروش فقط حضوری در یک مغازهٔ کوچک
        \item سفارش‌ها تلفنی و پراکنده
        \item مشتری نمی‌داند چه محصولی موجود است
        \item مقایسهٔ محصولات برای مشتری سخت است
        \item پیگیری گارانتی و سفارش دستی است
    \end{itemize}
}

% ---------- اسلاید ۱۵ ----------
\slide{
    \slidetitle{چشم‌انداز هوم‌استور}
    \vspace{1cm}
    \begin{center}
        {\Large\bfseries یک فروشگاه آنلاین لوازم خانگی}
        
        {\Large\bfseries با گارانتی و مقایسهٔ محصولات}
        
        \vspace{2cm}
        
        {\large مشتری می‌بیند ← مقایسه می‌کند ← سفارش می‌دهد}
        
        \vspace{0.5cm}
        
        {\large ادمین مدیریت می‌کند}
    \end{center}
}

% ---------- اسلاید ۱۶ ----------
\slide{
    \slidetitle{الهام از نمونه‌های واقعی}
    \vspace{1cm}
    \begin{itemize}[itemsep=0.8em]
        \item {\Large دیجی‌کالا}
        \item {\Large بامیلو}
        \item {\Large فروشگاه‌های لوازم خانگی آنلاین}
    \end{itemize}
    
    \vspace{1.5cm}
    
    \begin{center}
        {\Large\bfseries\myaccent{الهام بگیریم، کپی نکنیم.}}
    \end{center}
}

% ---------- اسلاید ۱۷ ----------
\slide{
    \slidetitle{محدودهٔ MVP هوم‌استور}
    \vspace{0.5cm}
    \begin{minipage}[t]{0.48\textwidth}
        \begin{minipage}[t]{0.48\textwidth}
            {\large\bfseries\textcolor{primary}{داخل محدوده:}}
            \vspace{0.3cm}
            \begin{itemize}[itemsep=0.3em]
                \item نمایش محصولات
                \item جست‌وجو و فیلتر
                \item مقایسهٔ محصولات
                \item سبد خرید
                \item پنل ادمین پایه
            \end{itemize}
        \end{minipage}
        \begin{minipage}[t]{0.48\textwidth}
            {\large\bfseries\myaccent{خارج از محدوده:}}
            \vspace{0.3cm}
            \begin{itemize}[itemsep=0.3em]
                \item پرداخت واقعی
                \item اپلیکیشن موبایل
                \item سیستم حمل و نقل
                \item گارانتی واقعی
            \end{itemize}
        \end{minipage}
    \end{minipage}
}

% ---------- اسلاید ۱۸ ----------
\slide{
    \slidetitle{تکنولوژی پیشنهادی}
    \vspace{0.5cm}
    \begin{itemize}[itemsep=0.6em]
        \item \textbf{Backend:} \lr{Node.js/Express} یا \lr{Python/Django}
        \item \textbf{Frontend:} \lr{React} یا \lr{Next.js}
        \item \textbf{DB:} \lr{PostgreSQL}
        \item \textbf{ابزار:} \lr{Git, GitHub, Figma, draw.io}
    \end{itemize}
    
    \vspace{1.5cm}
    
    \begin{center}
        {\large\bfseries\myaccent{آزادی انتخاب دارید، اما دلیل انتخاب را مستند کنید.}}
    \end{center}
}

% ---------- اسلاید ۱۹ ----------
\slide{
    \slidetitle{نقشهٔ سه Sprint هوم‌استور}
    \vspace{0.5cm}
    \begin{center}
        \begin{tabular}{clp{6cm}}
            \toprule
            \textbf{Sprint} & \textbf{هفته} & \textbf{هدف} \\
            \midrule
            1 & ۱–۵ & MVP: کاتالوگ + جست‌وجو + سبد خرید \\
            \hline
            2 & ۶–۱۱ & ثبت سفارش + مقایسه + پنل ادمین \\
            \hline
            3 & ۱۲–۱۴ & تست + استقرار + مستندات نهایی \\
            \bottomrule
        \end{tabular}
    \end{center}
}

% ---------- اسلاید ۲۰ ----------
\slide{
    \slidetitle{پروژهٔ شما: چه چیزی می‌سازید؟}
    \vspace{0.3cm}
    \begin{center}
        {\bfseries هر تیم ۳-۴ نفره، یک پروژهٔ مستقل انتخاب می‌کند}
    \end{center}
    
    \vspace{0.5cm}
    
    \begin{minipage}[t]{0.48\textwidth}
        \begin{minipage}[t]{0.48\textwidth}
            {\bfseries\textcolor{primary}{معیارهای خوب:}}
            \begin{itemize}[itemsep=0.3em]
                \item مشتری واقعی/فرضی
                \item حداقل ۳ نقش کاربری
                \item قابل تقسیم به ۳ Sprint
                \item حل مشکل واقعی
            \end{itemize}
        \end{minipage}
        \begin{minipage}[t]{0.48\textwidth}
            {\bfseries\myaccent{ایده‌های بد:}}
            \begin{itemize}[itemsep=0.3em]
                \item شبکهٔ اجتماعی
                \item بازی موبایل
                \item سیستم‌عامل جدید
                \item هر چیز طولانی
            \end{itemize}
        \end{minipage}
    \end{minipage}
}

% ---------- اسلاید ۲۱ ----------
\slide{
    \slidetitle{نمونه پروژه‌های خوب}
    \vspace{0.5cm}
    \begin{itemize}[itemsep=0.5em]
        \item سیستم رزرو نوبت آرایشگاه
        \item پلتفرم اجارهٔ کتاب
        \item اپ مدیریت وظایف تیمی
        \item سیستم سفارش آنلاین کافهٔ محلی
        \item فروشگاه صنایع دستی
        \item سیستم رزرو اتاق مطالعه
        \item اپ مدیریت اردوی دانشجویی
    \end{itemize}
}

% ---------- اسلاید ۲۲ ----------
\slide{
    \slidetitle{ترکیب نمره}
    \vspace{0.5cm}
    \begin{center}
        \begin{tabular}{cl}
            \toprule
            \textbf{مؤلفه} & \textbf{وزن} \\
            \midrule
            پروژه (۳ Sprint) & ۴۰٪ \\
            \hline
            کارگاه‌ها & ۱۵٪ \\
            \hline
            مشارکت و Peer Review & ۱۵٪ \\
            \hline
            میان‌ترم & ۱۵٪ \\
            \hline
            ارائهٔ نهایی & ۱۵٪ \\
            \bottomrule
        \end{tabular}
    \end{center}
}

% ---------- اسلاید ۲۳ ----------
\slide{
    \slidetitle{Definition of Done}
    \vspace{0.5cm}
    {\large یک کار وقتی «تمام» است که:}
    \vspace{0.5cm}
    \begin{itemize}[itemsep=0.5em]
        \item کد نوشته شده و Merge شده
        \item تست دارد و پاس می‌شود
        \item Code Review شده
        \item مستندات به‌روز است
        \item در محیط Demo کار می‌کند
    \end{itemize}
}

% ---------- اسلاید ۲۴ ----------
\slide{
    \slidetitle{نقش‌ها در تیم}
    \vspace{0.5cm}
    \begin{itemize}[itemsep=0.6em]
        \item \textbf{Product Owner:} تصمیم می‌گیرد چه چیزی بسازیم
        \item \textbf{Scrum Master:} جلسات و موانع را مدیریت می‌کند
        \item \textbf{Developer:} کد می‌نویسد
        \item \textbf{Tester:} کیفیت را تضمین می‌کند
    \end{itemize}
    
    \vspace{1.5cm}
    
    \begin{center}
        {\bfseries نقش‌ها هر Sprint می‌چرخند.}
    \end{center}
}

% ---------- اسلاید ۲۵ ----------
\slide{
    \slidetitle{قوانین همکاری}
    \vspace{0.5cm}
    \begin{itemize}[itemsep=0.6em]
        \item تأخیر در تحویل = تأثیر روی نمرهٔ تیم
        \item غیبت مکرر = حذف از تیم (موارد شدید)
        \item سهم نابرابر = گزارش به استاد
        \item استفاده از AI: مجاز، با شفافیت
        \item کپی از تیم دیگر = نمرهٔ صفر
    \end{itemize}
}

% ---------- اسلاید ۲۶ ----------
\slide{
    \slidetitle{ابزارها}
    \vspace{0.5cm}
    \begin{itemize}[itemsep=0.6em]
        \item \textbf{GitHub + Projects} — مدیریت کد و Backlog
        \item \textbf{Slack/Teams/Discord} — ارتباط تیمی
        \item \textbf{Figma} — طراحی UI
        \item \textbf{draw.io} — دیاگرام
        \item \textbf{Google Docs} — مستندات موقت
    \end{itemize}
}

% ---------- اسلاید ۲۷ ----------
\slide{
    \slidetitle{فعالیت کلاسی: بازی نقش}
    \vspace{1cm}
    \begin{center}
        {\Large\bfseries\myaccent{مرحلهٔ ۱ (۵ دقیقه):}}
        
        {\large من نقش خانم رضایی را بازی می‌کنم.}
        
        \vspace{1.5cm}
        
        {\Large\bfseries\myaccent{مرحلهٔ ۲ (۱۵ دقیقه):}}
        
        {\large هر تیم مشتری خودش را مصاحبه می‌کند.}
        
        \vspace{1.5cm}
        
        {\Large\bfseries هدف: کشف ۵ نیازمندی اصلی مشتری}
    \end{center}
}

% ---------- اسلاید ۲۸ ----------
\slide{
    \slidetitle{سؤالاتی که باید بپرسید}
    \vspace{0.5cm}
    \begin{itemize}[itemsep=0.7em]
        \item چه کسی مشتری اصلی شماست؟
        \item مشتری چه کاری می‌خواهد انجام دهد؟
        \item الان این کار را چطور انجام می‌دهد؟
        \item چه چیزی مهم‌تر است: سرعت، قیمت، یا کیفیت؟
        \item چه چیزی هرگز نباید اتفاق بیفتد؟
    \end{itemize}
}

% ---------- اسلاید ۲۹ ----------
\slide{
    \begin{center}
        \vspace{2.5cm}
        {\Huge\bfseries ما همین الان اولین قدم را برداشتیم.}
        
        \vspace{1.5cm}
        
        {\LARGE\myaccent{نیازمندی‌ها را کشف می‌کنیم، نه اختراع.}}
    \end{center}
}

% ---------- اسلاید ۳۰ ----------
\slide{
    \slidetitle{سه برداشت کلیدی}
    \vspace{1cm}
    \begin{enumerate}[itemsep=0.8em]
        \item {\large مهندسی نرم‌افزار = حل مسئله، نه فقط کد زدن.}
        \item {\large من پروژهٔ هوم‌استور را می‌سازم؛ شما پروژهٔ خودتان را.}
        \item {\large تیم شما مسئول موفقیت پروژهٔ خودش است.}
    \end{enumerate}
}

% ---------- اسلاید ۳۱ ----------
\slide{
    \slidetitle{تسک هفتهٔ آینده}
    \vspace{0.5cm}
    \begin{minipage}[t]{0.48\textwidth}
        \begin{minipage}[t]{0.48\textwidth}
            {\bfseries\textcolor{primary}{تسک شما:}}
            \vspace{0.3cm}
            \begin{itemize}[itemsep=0.3em]
                \item تیم ۳-۴ نفره
                \item مشتری مشخص
                \item Repository در GitHub
                \item \lr{01-vision.md}
                \item \lr{personas.md}
                \item ۱۰ User Story
            \end{itemize}
        \end{minipage}
        \begin{minipage}[t]{0.48\textwidth}
            {\bfseries\myaccent{تسک من:}}
            \vspace{0.3cm}
            \begin{itemize}[itemsep=0.3em]
                \item Repository هوم‌استور
                \item ساختار مستندات
                \item نسخهٔ اولیه در هفتهٔ ۲
                \item قالب Sprint Report
                \item نمونهٔ User Story
            \end{itemize}
        \end{minipage}
    \end{minipage}
}

% ---------- اسلاید ۳۲ ----------
\slide{
    \begin{center}
        \vspace{3cm}
        {\Huge\bfseries سؤالی دارید؟}
        
        \vspace{1.5cm}
        
        {\LARGE\myaccent{به امید یک ترم پربار}}
    \end{center}
}

\end{document}
